\section*{Задачи}
\addcontentsline{toc}{section}{Задачи}

\begin{enumerate}

    \item \textit{(Демо 2026, №9)} 
    Найти собственные векторы и собственные значения линейного оператора, заданного в некотором базисе матрицей
    \[
    A=\begin{pmatrix}
    0 & 2 \\
    -3 & 5
    \end{pmatrix}.
    \]
    Можно ли привести её к диагональному виду, перейдя к подходящему базису? Вычислить матрицу $A^n$, $n\in\mathbb{N}$.

    \item \textit{(Демо 2026, №10)}
    Найти матрицу линейного оператора, переводящего векторы
    \[
    a_1=(2,5)^T,\qquad a_2=(1,3)^T
    \]
    соответственно в векторы
    \[
    b_1=(7,-4)^T,\qquad b_2=(2,-1)^T
    \]
    в базисе, в котором даны координаты векторов.

    \item \textit{(Демо 2026, №11)}
    В базисе
    \[
    e_1=\begin{pmatrix} 3 \\ 1 \end{pmatrix},\qquad
    e_2=\begin{pmatrix} 2 \\ -1 \end{pmatrix}
    \]
    линейный оператор $\phi$ имеет матрицу
    \[
    A=\begin{pmatrix}
    -1 & 1 \\
    -3 & 4
    \end{pmatrix}.
    \]
    Найти матрицу оператора $\phi$ в базисе
    \[
    \hat e_1=\begin{pmatrix} 4 \\ 3 \end{pmatrix},\qquad
    \hat e_2=\begin{pmatrix} 1 \\ 1 \end{pmatrix}.
    \]

    \item \textit{(Демо 2026, №12)}
    Можно ли найти базис из собственных векторов для матрицы
    \[
    A=\begin{pmatrix}
    2 & -1 & 1\\
    -1 & 2 & -1\\
    1 & -1 & 2
    \end{pmatrix}?
    \]
    В случае положительного ответа найти этот базис, в случае отрицательного --- объяснить, почему это невозможно.

    \item \textit{(Демо 2026, №13)}
    Линейный оператор переводит вектор
    \[
    a_1=(-1,0)^T
    \]
    в вектор
    \[
    b_1=(5,5)^T,
    \]
    а вектор
    \[
    a_2=(1,1)^T
    \]
    в вектор
    \[
    b_2=(-2,-3)^T.
    \]
    \begin{enumerate}
        \item[1)] В какое множество перейдёт прямая, заданная уравнением
        \[
        2x_1-x_2=-2?
        \]
        \item[2)] Какое множество переходит в эту прямую?
        \item[3)] Написать уравнения тех прямых, которые переходят сами в себя.
    \end{enumerate}

    \item \textit{(Демо 2026, №14)}
    Найти базис ядра и базис образа линейного отображения
    \[
    \phi:\mathbb{R}^5\to\mathbb{R}^3,
    \]
    заданного матрицей
    \[
    A_\phi=
    \begin{pmatrix}
    1 & -1 & 2 & 4 & -2 \\
    3 & 9 & -14 & 2 & 1 \\
    3 & 6 & -9 & 1 & 1
    \end{pmatrix}.
    \]
    Является ли отображение сюръективным?

    \item \textit{(Демо 2026, №17)}
    Пусть в некотором ортонормированном базисе трёхмерного евклидова пространства заданы векторы
    \[
    e_1=(0,1,1)^T,\qquad e_2=(-1,-1,1)^T,\qquad e_3=(1,0,1)^T.
    \]
    Пусть оператор $f$ задан матрицей
    \[
    A_f=
    \begin{pmatrix}
    3 & 1 & 1 \\
    1 & 5 & 0 \\
    -3 & 2 & 7
    \end{pmatrix}
    \]
    в базисе $e_1,e_2,e_3$.
    Найти матрицу $A_{f^*}$ сопряжённого оператора $f^*$ в том же базисе.
\end{enumerate}
